\documentclass{article}
\usepackage{amsmath,amssymb,xurl}
\usepackage[hidelinks]{hyperref}
% Shared commands for notes in docs/paper-gaps/.

\DeclareUnicodeCharacter{2080}{\ifmmode{}_{0}\else\textsubscript{0}\fi}
\DeclareUnicodeCharacter{2081}{\ifmmode{}_{1}\else\textsubscript{1}\fi}
\DeclareUnicodeCharacter{2082}{\ifmmode{}_{2}\else\textsubscript{2}\fi}
\DeclareUnicodeCharacter{2099}{\ifmmode{}_{n}\else\textsubscript{n}\fi}

\newcommand{\Lean}{\textsc{Lean}}
\ExplSyntaxOn
\NewDocumentCommand{\leanid}{m}{
  \ifmmode
    \textsf{\detokenize{#1}}
  \else
    \group_begin:
    \urlstyle{sf}
    \tl_set:Nn \l_tmpa_tl {#1}
    \tl_replace_all:Nnn \l_tmpa_tl {\_} {_}
    \exp_args:NV \path \l_tmpa_tl
    \group_end:
  \fi
}
\ExplSyntaxOff
\providecommand{\tr}{\operatorname{tr}}
\newcommand{\tnleanIssueBase}{https://github.com/LionSR/TNLean/issues/}
\newcommand{\tnleanPrBase}{https://github.com/LionSR/TNLean/pull/}
\newcommand{\tnleanDocBase}{https://sirui-lu.com/TNLean/docs/}
\newcommand{\ghissue}[1]{\href{\tnleanIssueBase#1}{GitHub issue \##1}}
\newcommand{\ghpr}[1]{\href{\tnleanPrBase#1}{PR \##1}}
\newcommand{\doclink}[1]{\href{\tnleanDocBase#1.html}{\path{#1}}}

% Dirac notation and the identity operator, as in blueprint/src/macros/common.tex.
% \providecommand keeps note-local definitions authoritative.
\usepackage{dsfont}
\providecommand{\ket}[1]{|#1\rangle}
\providecommand{\bra}[1]{\langle #1|}
\providecommand{\braket}[2]{\langle #1 | #2 \rangle}
\providecommand{\ketbra}[2]{|#1\rangle\!\langle #2|}
\providecommand{\Id}{\mathds{1}}
\providecommand{\one}{\mathds{1}}
\providecommand{\C}{\mathbb{C}}
\providecommand{\MN}[1]{M_{#1}(\C)}
\providecommand{\MD}{M_D(\C)}

% Verdict marker: \gapnote{<kind>}{<status>}. Kind is the policy
% classification (clarification, local-correction, scope-restriction,
% unfaithful, false-source, open-gap); status is open, wip (elimination
% underway), resolved, or historical. Machine-read by the site index;
% prints a verdict line.
\newcommand{\gapnote}[2]{\par\noindent\textbf{Verdict:} #1 (#2).\par}

\title{Locally Flat Transports and Restricted Dual Homotopy}
\author{}
\date{}
\begin{document}
\maketitle
\gapnote{scope-restriction}{open}

Schuch, Cirac, and P\'erez-Garc\'ia,
\href{https://arxiv.org/abs/1001.3807v3}{arXiv:1001.3807v3},
Lemma 6.14, source \path{Papers/1001.3807/paper_v3.tex}, lines 2199--2214,
deforms flux strings using virtual invariance. When the swept support is
restricted, equal lifted displacement alone does not imply the generated
homotopy: a missing square can retain a nonabelian loop obstruction.

The definition \leanid{TNLean.PEPS.TorusDualFlatTransport} instead specifies
independent local data: a set at each dual vertex, mutually inverse maps
on opposite labelled steps, and equal compositions on the two sides of
every allowed square. These conditions make no reference to a pair of
paths or a homotopy witness. The definition
\leanid{TNLean.PEPS.TorusDualFlatTransport.eval} composes the step maps.
The theorem \leanid{TNLean.PEPS.TorusDualHomotopy.flatTransport_eval_eq}
derives transport equality from every generated supported homotopy.

The converse, \leanid{TNLean.PEPS.torusDualHomotopy_iff_flatTransport},
proves positive semantic completeness under all such locally flat
transports. Its separating transport is constructed from the classes of
actual paths from the common initial vertex, with the existing generated
homotopy as equivalence relation. Appending a step is reversible by an
elementary backtrack and locally flat by an allowed square. Evaluating at
the constant path class recovers a homotopy of the compared paths.
This construction uses no torus size assumption, including for the
separately labelled parallel steps of periods one and two.

The universal family includes arbitrary vertex sets, rather than only
finite groups, a finite collection of seam words, or physical flux
insertions. Thus this theorem supplies an exact transport criterion, not
a finite decision procedure or a claim that winding is complete under
restricted support. It does not derive a prescribed geometric homotopy
from endpoint exclusion alone and does not prove noncentral open-pair
vacuum density. Those source-level obligations remain under
\ghissue{8703}; the full-lattice Hamiltonian scope of \ghissue{8947} is
unchanged.
\end{document}
